\section{Discussion}
\label{sec:discussion}

\subsection{Lessons}

\paragraph{Report detection per class, and per device}
%%LESSON1%%

\paragraph{Evaluate the running system, not only the model}
A model can be evaluated flawlessly on its dataset while the deployed system
feeds it something else. \textsc{Sentinel}'s extractor was written from a
description of Kitsune rather than from its code, and the error it contained
was invisible in every offline experiment, because offline experiments read
N-BaIoT's precomputed features and never touch the extractor. The check that
caught it needs no additional data: structural invariants measured on the
training set's own rows, asserted on live output. We recommend such parity
tests as a standing part of any pipeline that trains on features computed by
someone else's tool.

\paragraph{Recover what an artefact encodes}
The deployed model arrived without its training script. Its scaler recorded,
in its per-feature minima, maxima and sample count, enough to identify the
exact training devices---which in turn exposed that its in-distribution
evaluation would have been contaminated, and supplied a natural held-out-device
split. Fitted preprocessing objects are an underused source of provenance.

\paragraph{Give the model the evidence}
%%LESSON4%%

\paragraph{Latency lives in the integration}
Every component we wrote runs in well under a millisecond per event; the
detection-to-response path is slow because of an external reputation service
called in the wrong way and a language model on a CPU. Neither problem is
visible when the detector is benchmarked alone. Requesting only the needed
section of the OTX response, running it off the event loop with a timeout, and
caching per address would remove most of the reputation cost.

\subsection{Limitations and threats to validity}
\emph{Dataset.} N-BaIoT was captured in 2018 in a laboratory, with nine
devices and two botnet families; home networks today carry more devices, more
encrypted traffic and other malware. Our results characterise these attacks on
these devices. The held-out-device split measures generalisation to three
unseen devices of the same era, not to arbitrary new hardware.

\emph{Hardware.} All timings come from a four-core aarch64 virtual machine,
not a Raspberry~Pi~5. The host matches the Pi in instruction set, core count
and approximate memory, but not in clock speed, memory bandwidth or thermal
behaviour, and sustained LLM inference on the Pi will be slower. Physical
validation remains future work.

\emph{Replay, not deployment.} Detection results come from replaying captured
features, not from a live network; the parity test narrows but does not close
the gap between the two. We have not evaluated \textsc{Sentinel} against an
adaptive attacker.

\emph{LLM evaluation.} The ground truth maps attack classes to ATT\&CK
techniques by our reading of the dataset's definitions; other mappings are
defensible, which is why we accept several techniques per class and report
tactic-level accuracy as well. The model samples at temperature~1.0, so a
single run estimates, rather than fixes, its behaviour; 200 windows per prompt
give 95\% intervals no wider than about $\pm$7 percentage points.

\emph{Deception path.} The honeypot-to-firewall path was exercised
functionally, by attacking our own honeypot, but is not evaluated
quantitatively here; the appliance's own event log mixes those sessions with
synthetic events injected to test the response logic, so we do not report
counts from it.

\emph{Single threshold.} Every retrained model uses one global threshold at a
nominal 1\% false-positive budget; other operating points trade detection for
false positives differently, and the threshold-free AUC values we report
summarise that trade-off.

\subsection{Ethics}
No data about people was collected for this work. Detection experiments use
the public N-BaIoT dataset and the Mirai capture distributed with Kitsune. The
LLM experiment used reserved documentation addresses (RFC~5737), and
reputation lookups in the latency experiment were made only for such
addresses. The honeypot path was exercised only by the author, attacking the
author's own honeypot from the author's own device.
